% 结论与政策建议
\section{结论与政策建议}

\subsection{主要发现}

本文利用 2011--2023 年中国 31 个省份平衡面板数据，采用双向固定效应模型和 Hansen（1999）面板门槛模型，实证检验了数字经济发展对省际就业结构的影响。主要发现可概括为以下四点。

第一，数字经济发展显著提升了省际第三产业就业占比，验证了 H1。双向固定效应基准模型（M6）的估计结果表明，数字经济综合指数每提高 1 个单位，三产就业占比提高约 0.277 个单位（$p<0.01$）。一个标准差（0.123）的数字经济指数增长对应三产就业占比提升约 3.4 个百分点。该结果在 M1$\rightarrow$M6 六种渐进控制规格中均保持方向一致和统计显著，且经替代缩尾水平、对数变换和剔除 COVID 年份三项稳健性检验后保持稳健。

第二，数字经济的就业结构效应存在显著的区域梯度格局，部分支持了 H3。东部地区的边际效应为 0.304（$p<0.05$），与全样本基准接近；中部地区的边际效应高达 1.781（$p<0.05$），系数量级约为东部的六倍，提示中部地区正处于数字化驱动的就业结构快速转换期；西部地区系数为 0.158，不显著。这一东--中--西梯度格局与技术扩散理论的"追赶者红利"假说一致——技术溢出效应在技术基础较好但尚未达到前沿的地区最为显著。

第三，广义数字经济综合指数对就业结构的影响显著强于狭义数字金融指数。替换为核心解释变量后，数字金融指数的系数仅为 0.055 且不显著（$p=0.153$）。这一对比揭示了一个政策相关洞见：推动就业结构转型的数字化驱动力更多来自互联网基础设施、移动通信覆盖和软件服务等广义数字化维度，而非单纯的金融科技推广。

第四，城镇化水平对数字经济--就业结构关系存在调节效应，但统计显著性未达标，H2 仅获部分支持。门槛分析的最优门槛值为城镇化率 82.3\%，两端系数差异与预期方向一致（高城镇化组 0.388 > 低城镇化组 0.200），但高组仅包含 3 个直辖市（35 个观测值），Bootstrap 门槛显著性检验因聚类数过少未能收敛。这一结果更宜被视为一种启发性的模式识别而非严格的因果门槛推断。

\subsection{政策含义}

上述发现为政策制定者提供了三个维度的参考框架。

第一，数字经济发展应成为促进就业结构转型升级的重要政策工具。双向固定效应下的稳健正向系数意味着，即使在控制了经济发展水平、资本形成、开放度、政府规模和人力资本投资后，数字经济指数与三产就业占比之间仍存在显著的条件性正相关。对于面临传统制造业收缩、就业结构转型压力的省份，加大数字经济基础设施建设可能产生正向的就业外溢效应。

第二，"抓广度优于抓深度"——广义数字基础设施（宽带、移动通信、软件产业生态）对就业结构的推动作用大于单一的金融科技应用。这一发现对地方政府制定数字经济战略具有直接的政策优先序含义：在有限的公共财政资源约束下，优先投资于 5G 基站、千兆光网、数据中心等通用数字基础设施，可能比押注特定的数字金融应用推广更有效地带动就业结构升级。

第三，区域差异化政策设计值得重视。中部地区（湖北、湖南、河南、安徽、江西、山西等）正处于数字化的"甜蜜点"——基础设施已基本到位、人力资本正在追赶、产业数字化的需求旺盛。在这些省份加大数字经济投入，可以以较少的增量投资换取较大的就业结构转换效应。西部地区则可能需要先完成基础性的数字接入（主要是宽带和移动通信的覆盖率和可负担性），然后才能期待数字经济发挥显著的就业效益。

\subsection{研究局限与未来展望}

本文存在四个主要局限，为后续研究指明了方向。

第一，被解释变量的维度受限。由于数据可及性限制，本文仅检验了就业结构（第三产业就业占比）维度，未能涵盖就业规模（总就业人数或城镇就业人数）和就业质量（失业率、工资水平、就业稳定性等）两个传统劳动经济学核心维度。这一限制使得本文无法回答"数字经济发展对总体就业规模是净创造还是净消灭"这一政策界最关心的问题。后续研究如能获取更全面的省际就业变量，将有助于形成一个更完整的数字就业效应画像。一个可行的路径是使用省级人口普查或 1\% 人口抽样调查的微观数据，将本文的省级 FE 框架扩展为县级或个体层面的因果识别设计。

第二，内生性处理不充分。本文的识别策略依赖双向固定效应下的条件均值独立假设——即控制了省份 FE、年份 FE 和八个控制变量后，残余的数字经济变异与残余的就业结构变异条件独立。这一假设无法直接检验。两个具体的威胁值得关注：反向因果（就业结构先进 → 服务业主导 → 更多数字技术需求 → 数字经济发展）和时变遗漏变量（如省内市场化程度改革、优惠产业政策等同时对数字经济和就业结构产生影响的第三因素）。后续研究可考虑采用工具变量策略，一个可行的思路是利用全国数字经济指数的年度增长率作为各省指数增长的工具变量——全国趋势反映技术前沿的普遍进步，不太可能由单个省份的就业结构变化驱动——或者利用 Bartik（shift-share）设计，以各省初始产业数字化的强度与全国行业数字化增长的交互作为工具变量。

第三，门槛分析的统计效力和样本覆盖不足。Hansen 门槛模型需要在两个区间内分别有充分的聚类数和观测值来保证估计和推断的可靠性。本文的城镇化门槛 82.3\% 仅划分出 3 个省份作为高城镇化组，使得 Bootstrap 检验无法完成。一个改进路径是使用地级市层面数据（$N\approx 3000$），将获得足够的高城镇化城市样本来进行可信的门槛推断。另一个路径是考虑放弃离散门槛设定，转而使用半参数或非参数方法（如局部线性核回归或交互固定效应模型）来估计城镇化对数字经济效应的连续调节函数。

第四，本文以省际宏观面板为分析单元，这意味着数字经济就业效应的微观机制无法被直接检验。宏观层面的正相关可能由多个微观渠道共同驱动：数字平台降低了劳动力市场的信息不对称（匹配效率提升）、数字技能的溢价改变了教育投资的回报率（技能结构转变）、共享经济和零工经济改变了就业形态的统计归属（统计口径变化）。要区分这些机制，需要将省际宏观面板分析延伸到地级市层面甚至企业层面，并利用特定的自然实验（如"宽带中国"试点政策、数字经济综合试验区设立等准实验）进行因果机制的分解检验。
